\chapter{El mapa del tag y las máscaras}\label{cap:mapa}

Este capítulo es el pegamento entre las estaciones. La librería solo transporta bytes; aquí se explica qué significa cada uno, cómo elige cada estación los suyos, y cómo se calculan las máscaras.

\section{El bloque de 28 bytes}
El grupo acordó (hoja de cálculo \texttt{Mapa\_RFID\_CPLab.xlsx}) usar un único bloque de 28 bytes en el tag. Los bytes 0 a 19 son del grupo; los bytes 20 a 27 están reservados para compartir con los otros grupos de la CP-Lab. La figura~\ref{fig:bloque} los muestra con un color por quién los usa.

\begin{figure}[ht]
\centering
\begin{tikzpicture}[x=0.585cm,y=1cm]
\fill[gris!35] (0,-0.00) rectangle (1,0.80);\draw[gray!70] (0,-0.00) rectangle (1,0.80);
\node[font=\scriptsize\bfseries] at (0.5,0.40) {0};
\fill[gris!35] (1,-0.00) rectangle (2,0.80);\draw[gray!70] (1,-0.00) rectangle (2,0.80);
\node[font=\scriptsize\bfseries] at (1.5,0.40) {1};
\fill[gris!35] (2,-0.00) rectangle (3,0.80);\draw[gray!70] (2,-0.00) rectangle (3,0.80);
\node[font=\scriptsize\bfseries] at (2.5,0.40) {2};
\fill[gris!35] (3,-0.00) rectangle (4,0.80);\draw[gray!70] (3,-0.00) rectangle (4,0.80);
\node[font=\scriptsize\bfseries] at (3.5,0.40) {3};
\fill[gris!35] (4,-0.00) rectangle (5,0.80);\draw[gray!70] (4,-0.00) rectangle (5,0.80);
\node[font=\scriptsize\bfseries] at (4.5,0.40) {4};
\fill[gris!35] (5,-0.00) rectangle (6,0.80);\draw[gray!70] (5,-0.00) rectangle (6,0.80);
\node[font=\scriptsize\bfseries] at (5.5,0.40) {5};
\fill[gris!35] (6,-0.00) rectangle (7,0.80);\draw[gray!70] (6,-0.00) rectangle (7,0.80);
\node[font=\scriptsize\bfseries] at (6.5,0.40) {6};
\fill[gris!35] (7,-0.00) rectangle (8,0.80);\draw[gray!70] (7,-0.00) rectangle (8,0.80);
\node[font=\scriptsize\bfseries] at (7.5,0.40) {7};
\fill[cam!35] (8,-0.00) rectangle (9,0.80);\draw[gray!70] (8,-0.00) rectangle (9,0.80);
\node[font=\scriptsize\bfseries] at (8.5,0.40) {8};
\fill[cam!35] (9,-0.00) rectangle (10,0.80);\draw[gray!70] (9,-0.00) rectangle (10,0.80);
\node[font=\scriptsize\bfseries] at (9.5,0.40) {9};
\fill[cam!35] (10,-0.00) rectangle (11,0.80);\draw[gray!70] (10,-0.00) rectangle (11,0.80);
\node[font=\scriptsize\bfseries] at (10.5,0.40) {10};
\fill[cam!35] (11,-0.00) rectangle (12,0.80);\draw[gray!70] (11,-0.00) rectangle (12,0.80);
\node[font=\scriptsize\bfseries] at (11.5,0.40) {11};
\fill[back!35] (12,-0.00) rectangle (13,0.80);\draw[gray!70] (12,-0.00) rectangle (13,0.80);
\node[font=\scriptsize\bfseries] at (12.5,0.40) {12};
\fill[press!35] (13,-0.00) rectangle (14,0.80);\draw[gray!70] (13,-0.00) rectangle (14,0.80);
\node[font=\scriptsize\bfseries] at (13.5,0.40) {13};
\fill[press!35] (0,-1.55) rectangle (1,-0.75);\draw[gray!70] (0,-1.55) rectangle (1,-0.75);
\node[font=\scriptsize\bfseries] at (0.5,-1.15) {14};
\fill[press!35] (1,-1.55) rectangle (2,-0.75);\draw[gray!70] (1,-1.55) rectangle (2,-0.75);
\node[font=\scriptsize\bfseries] at (1.5,-1.15) {15};
\fill[press!35] (2,-1.55) rectangle (3,-0.75);\draw[gray!70] (2,-1.55) rectangle (3,-0.75);
\node[font=\scriptsize\bfseries] at (2.5,-1.15) {16};
\fill[out!35] (3,-1.55) rectangle (4,-0.75);\draw[gray!70] (3,-1.55) rectangle (4,-0.75);
\node[font=\scriptsize\bfseries] at (3.5,-1.15) {17};
\fill[out!35] (4,-1.55) rectangle (5,-0.75);\draw[gray!70] (4,-1.55) rectangle (5,-0.75);
\node[font=\scriptsize\bfseries] at (4.5,-1.15) {18};
\fill[gris!35] (5,-1.55) rectangle (6,-0.75);\draw[gray!70] (5,-1.55) rectangle (6,-0.75);
\node[font=\scriptsize\bfseries] at (5.5,-1.15) {19};
\fill[comp!35] (6,-1.55) rectangle (7,-0.75);\draw[gray!70] (6,-1.55) rectangle (7,-0.75);
\node[font=\scriptsize\bfseries] at (6.5,-1.15) {20};
\fill[comp!35] (7,-1.55) rectangle (8,-0.75);\draw[gray!70] (7,-1.55) rectangle (8,-0.75);
\node[font=\scriptsize\bfseries] at (7.5,-1.15) {21};
\fill[comp!35] (8,-1.55) rectangle (9,-0.75);\draw[gray!70] (8,-1.55) rectangle (9,-0.75);
\node[font=\scriptsize\bfseries] at (8.5,-1.15) {22};
\fill[comp!35] (9,-1.55) rectangle (10,-0.75);\draw[gray!70] (9,-1.55) rectangle (10,-0.75);
\node[font=\scriptsize\bfseries] at (9.5,-1.15) {23};
\fill[comp!35] (10,-1.55) rectangle (11,-0.75);\draw[gray!70] (10,-1.55) rectangle (11,-0.75);
\node[font=\scriptsize\bfseries] at (10.5,-1.15) {24};
\fill[comp!35] (11,-1.55) rectangle (12,-0.75);\draw[gray!70] (11,-1.55) rectangle (12,-0.75);
\node[font=\scriptsize\bfseries] at (11.5,-1.15) {25};
\fill[comp!35] (12,-1.55) rectangle (13,-0.75);\draw[gray!70] (12,-1.55) rectangle (13,-0.75);
\node[font=\scriptsize\bfseries] at (12.5,-1.15) {26};
\fill[comp!35] (13,-1.55) rectangle (14,-0.75);\draw[gray!70] (13,-1.55) rectangle (14,-0.75);
\node[font=\scriptsize\bfseries] at (13.5,-1.15) {27};
\end{tikzpicture}
\\[6pt]
\small
\tikz\fill[gris!35] (0,0) rectangle (0.35,0.25);\ General\quad
\tikz\fill[cam!35] (0,0) rectangle (0.35,0.25);\ Cámara\quad
\tikz\fill[back!35] (0,0) rectangle (0.35,0.25);\ BackCover\quad
\tikz\fill[press!35] (0,0) rectangle (0.35,0.25);\ Press\quad
\tikz\fill[out!35] (0,0) rectangle (0.35,0.25);\ Out\quad
\tikz\fill[comp!35] (0,0) rectangle (0.35,0.25);\ Compartido (otros grupos)
\caption{Los 28 bytes del bloque, byte a byte (0 a 13 arriba, 14 a 27 abajo).}
\label{fig:bloque}
\end{figure}

\begin{idea}
El tag es un \textbf{tablón compartido}. Cada byte tiene un dueño que lo escribe, y varios que lo leen. Que un byte tenga un solo dueño es lo que impide que las estaciones se pisen. La excepción son los bytes de «banderas» (el byte 4, que cada estación completa con su bit); para esos hay una regla especial (sección~\ref{sec:reglas-oro}).
\end{idea}

\section{El mapa completo}
La tabla siguiente es el contenido de la hoja «Mapa RFID» del Excel del grupo. Las filas están coloreadas por estación. \textbf{Si alguna vez la tabla de este documento y el Excel difieren, manda el Excel}, que es donde se acuerdan los cambios.

{\small\setlength{\tabcolsep}{3.2pt}
\begin{longtable}{>{\raggedleft\arraybackslash}p{0.9cm}p{3.6cm}p{2.0cm}p{2.5cm}p{2.3cm}p{3.5cm}}
\toprule
\textbf{Byte} & \textbf{Campo} & \textbf{Estación} & \textbf{Lo escribe} & \textbf{Lo lee} & \textbf{Valores} \\
\midrule
\endhead
\rowcolor{gris!18}0 & Marca de bloque válido & General & 1\textordfeminine{} estación que lo inicializa & Todas & 16\#A5 = bloque inicializado \\
\rowcolor{gris!18}1 & ID de palet / n\textordmasculine{} de pieza & General & HMI o 1\textordfeminine{} estación & Todas & 1–255 \\
\rowcolor{gris!18}2 & Receta / tipo de producto & General & HMI o 1\textordfeminine{} estación & Todas & Códigos a definir por el grupo \\
\rowcolor{gris!18}3 & Estaciones a realizar (máscara) & General & HMI o 1\textordfeminine{} estación & Todas & Ver hoja 'Máscaras' \\
\rowcolor{gris!18}4 & Estaciones completadas (máscara) & General & Cada estación pone su bit & Todas & Ver hoja 'Máscaras' \\
\rowcolor{gris!18}5 & Calidad global & General & Cámara & Out & 0 sin evaluar \,\textperiodcentered\, 1 OK \,\textperiodcentered\, 2 NOK \\
\rowcolor{gris!18}6 & Código de rechazo & General & Estación que rechaza & Out, HMI & Nibble alto = estación (1 Cám, 2 Back, 3 Press, 4 Out) \,\textperiodcentered\, nibble bajo = motivo \\
\rowcolor{gris!18}7 & Libre del grupo & General & — & — & Libre \\
\rowcolor{cam!18}8 & JobNumber de la cámara & Cámara & Receta / HMI & Cámara & 1 = 0 fusibles \,\textperiodcentered\, 2 = izq \,\textperiodcentered\, 3 = dcha \,\textperiodcentered\, 4 = 2 fusibles \,\textperiodcentered\, 5 = orientación palet \\
\rowcolor{cam!18}9 & Resultado de la cámara & Cámara & Cámara & Out & bit0 inspeccionado \,\textperiodcentered\, bit1 OK (RDBU) \,\textperiodcentered\, bit2 orientación OK \\
\rowcolor{cam!18}10 & ImageID (\%EB1004) & Cámara & Cámara & HMI / trazabilidad & 0–255 \\
\rowcolor{cam!18}11 & Intentos de inspección & Cámara & Cámara & HMI & 0–n \\
\rowcolor{back!18}12 & Estado BackCover & BackCover & BackCover & Press & bit0 requerida \,\textperiodcentered\, bit1 colocada \,\textperiodcentered\, bit2 falta stock \\
\rowcolor{press!18}13 & N\textordmasculine{} de golpes de prensa & Press & Receta / HMI & Press & 1–n \\
\rowcolor{press!18}14--15 & Tiempo de prensado (ms) & Press & Receta / HMI & Press & p. ej. 2000 \\
\rowcolor{press!18}16 & Golpes realizados & Press & Press & Out & 0–n \\
\rowcolor{out!18}17 & Almacén destino & Out & Receta (opcional) & Out & 0 auto por calidad \,\textperiodcentered\, 1 izquierdo \,\textperiodcentered\, 2 derecho \\
\rowcolor{out!18}18 & Almacén depositado & Out & Out & HMI & 1 izquierdo \,\textperiodcentered\, 2 derecho \\
\rowcolor{gris!18}19 & Libre del grupo & General & — & — & Libre \\
\rowcolor{comp!18}20 & RESERVADO – Fusibles colocados (propuesta) & Compartido & RASS (otro grupo) & Cámara & bit0 izquierdo \,\textperiodcentered\, bit1 derecho \\
\rowcolor{comp!18}21 & RESERVADO – Resultado medición (propuesta) & Compartido & Measurement (otro grupo) & Out & 0 sin medir \,\textperiodcentered\, 1 OK \,\textperiodcentered\, 2 NOK \\
\rowcolor{comp!18}22 & RESERVADO – N\textordmasculine{} de taladros (propuesta) & Compartido & Drilling (otro grupo) & Out (en CIROS) & 0, 2, 4 \\
\rowcolor{comp!18}23 & RESERVADO – Tapa frontal colocada (propuesta) & Compartido & FrontCover (otro grupo) & BackCover & 0 no \,\textperiodcentered\, 1 sí \\
\rowcolor{comp!18}24--27 & RESERVADO – Libre para compartir & Compartido & Por acordar & Por acordar & Libre \\
\bottomrule
\end{longtable}

}

\begin{ojo}
Los dos bytes 14--15 forman un único número (\texttt{INT}) en formato \emph{big-endian}: el byte 14 es la parte alta y el 15 la baja. 2000\,ms es \cod{16\#07D0}, es decir, byte 14 = \cod{07} y byte 15 = \cod{D0}. Léelos y escríbelos siempre juntos. Los campos de este tipo empiezan en un byte par para poder leerse de una vez.
\end{ojo}

\section{Qué hace cada estación con los bytes}
\begin{tabularx}{\textwidth}{>{\bfseries}l p{2.6cm} p{3.6cm} X}
\toprule
Estación & Lee (bytes) & Escribe (bytes) & Secuencia resumida \\
\midrule
Cámara & 0, 3, 8, 20 & 0–4 si el tag está vacío; 4 (bit0), 5, 6, 9, 10, 11 & Si el byte 0 no es 16\#A5, inicializar el bloque con valores por defecto. Leer el job (byte 8); si es 0, deducirlo de los fusibles (byte 20) o usar el job por defecto. Switch-to-run $\rightarrow$ Change Job $\rightarrow$ Trigger. Esperar xReady y leer RDBU. LED verde si OK, rojo parpadeante si NOK. Escribir resultado, calidad e ImageID. \\
\midrule
BackCover & 3, 23 & 4 (bit1), 12 & Comprobar que está requerida (byte 3 bit1), que hay tapa frontal (sensor; byte 23 si se acuerda) y que no hay trasera. Colocar la tapa. Si el almacén está vacío, poner el bit2 del byte 12. \\
\midrule
Press & 3, 12, 13, 14–15 & 4 (bit2), 16 & Comprobar que hay tapa trasera (byte 12 bit1 y sensor). Leer n\textordmasculine{} de golpes y tiempo (0 $\rightarrow$ valores por defecto). Subir, bajar hasta presión inferior, mantener, subir. Repetir n veces. Escribir golpes realizados. \\
\midrule
Out & 5, 9, 17, 21, 22 & 4 (bit3), 18; al final todo el bloque a 0 & Destino = byte 17; si es 0: calidad OK$\rightarrow$almacén izquierdo, NOK$\rightarrow$derecho (en CIROS: 2 vs 4 taladros, byte 22). Mover X solo con Z arriba. Tras depositar, escribir byte 18 y poner a cero el bloque. \\
\midrule
 &  &  &  \\
\midrule
 &  &  &  \\
\midrule
 &  &  &  \\
\midrule
 &  &  &  \\
\midrule
 &  &  &  \\
\midrule
 &  &  &  \\
\midrule
 &  &  &  \\
\midrule
 &  &  &  \\
\midrule
 &  &  &  \\
\midrule
 &  &  &  \\
\midrule
 &  &  &  \\
\midrule
 &  &  &  \\
\midrule
 &  &  &  \\
\midrule
 &  &  &  \\
\midrule
 &  &  &  \\
\midrule
 &  &  &  \\
\midrule
 &  &  &  \\
\midrule
 &  &  &  \\
\midrule
 &  &  &  \\
\midrule
 &  &  &  \\
\midrule
 &  &  &  \\
\midrule
 &  &  &  \\
\midrule
 &  &  &  \\
\midrule
 &  &  &  \\
\midrule
 &  &  &  \\
\midrule
 &  &  &  \\
\midrule
 &  &  &  \\
\midrule
 &  &  &  \\
\midrule
 &  &  &  \\
\midrule
 &  &  &  \\
\midrule
 &  &  &  \\
\midrule
 &  &  &  \\
\midrule
 &  &  &  \\
\midrule
 &  &  &  \\
\midrule
 &  &  &  \\
\midrule
 &  &  &  \\
\midrule
 &  &  &  \\
\midrule
 &  &  &  \\
\midrule
 &  &  &  \\
\midrule
 &  &  &  \\
\midrule
 &  &  &  \\
\midrule
 &  &  &  \\
\midrule
 &  &  &  \\
\midrule
 &  &  &  \\
\midrule
 &  &  &  \\
\midrule
 &  &  &  \\
\midrule
 &  &  &  \\
\midrule
 &  &  &  \\
\midrule
 &  &  &  \\
\midrule
 &  &  &  \\
\midrule
 &  &  &  \\
\midrule
 &  &  &  \\
\midrule
 &  &  &  \\
\midrule
 &  &  &  \\
\midrule
 &  &  &  \\
\midrule
 &  &  &  \\
\midrule
 &  &  &  \\
\midrule
 &  &  &  \\
\midrule
 &  &  &  \\
\midrule
 &  &  &  \\
\midrule
 &  &  &  \\
\midrule
 &  &  &  \\
\midrule
 &  &  &  \\
\midrule
 &  &  &  \\
\midrule
 &  &  &  \\
\midrule
 &  &  &  \\
\midrule
 &  &  &  \\
\midrule
 &  &  &  \\
\midrule
 &  &  &  \\
\midrule
 &  &  &  \\
\midrule
 &  &  &  \\
\midrule
 &  &  &  \\
\midrule
 &  &  &  \\
\midrule
 &  &  &  \\
\midrule
 &  &  &  \\
\midrule
 &  &  &  \\
\midrule
 &  &  &  \\
\midrule
 &  &  &  \\
\midrule
 &  &  &  \\
\midrule
 &  &  &  \\
\midrule
 &  &  &  \\
\midrule
 &  &  &  \\
\midrule
 &  &  &  \\
\midrule
 &  &  &  \\
\midrule
 &  &  &  \\
\midrule
 &  &  &  \\
\midrule
 &  &  &  \\
\midrule
 &  &  &  \\
\midrule
 &  &  &  \\
\midrule
 &  &  &  \\
\midrule
 &  &  &  \\
\midrule
 &  &  &  \\
\midrule
 &  &  &  \\
\midrule
 &  &  &  \\
\midrule
 &  &  &  \\
\midrule
 &  &  &  \\
\midrule
 &  &  &  \\
\midrule
 &  &  &  \\
\midrule
 &  &  &  \\
\midrule
 &  &  &  \\
\midrule
 &  &  &  \\
\midrule
 &  &  &  \\
\midrule
 &  &  &  \\
\midrule
 &  &  &  \\
\midrule
 &  &  &  \\
\midrule
 &  &  &  \\
\midrule
 &  &  &  \\
\midrule
 &  &  &  \\
\midrule
 &  &  &  \\
\midrule
 &  &  &  \\
\midrule
 &  &  &  \\
\midrule
 &  &  &  \\
\midrule
 &  &  &  \\
\midrule
 &  &  &  \\
\midrule
 &  &  &  \\
\midrule
 &  &  &  \\
\midrule
 &  &  &  \\
\midrule
 &  &  &  \\
\midrule
 &  &  &  \\
\midrule
 &  &  &  \\
\midrule
 &  &  &  \\
\midrule
 &  &  &  \\
\midrule
 &  &  &  \\
\midrule
 &  &  &  \\
\midrule
 &  &  &  \\
\midrule
 &  &  &  \\
\midrule
 &  &  &  \\
\midrule
 &  &  &  \\
\midrule
 &  &  &  \\
\midrule
 &  &  &  \\
\midrule
 &  &  &  \\
\midrule
 &  &  &  \\
\midrule
 &  &  &  \\
\midrule
 &  &  &  \\
\midrule
 &  &  &  \\
\midrule
 &  &  &  \\
\midrule
 &  &  &  \\
\midrule
 &  &  &  \\
\midrule
 &  &  &  \\
\midrule
 &  &  &  \\
\midrule
 &  &  &  \\
\midrule
 &  &  &  \\
\midrule
 &  &  &  \\
\midrule
 &  &  &  \\
\midrule
 &  &  &  \\
\midrule
 &  &  &  \\
\midrule
 &  &  &  \\
\midrule
 &  &  &  \\
\midrule
 &  &  &  \\
\midrule
 &  &  &  \\
\midrule
 &  &  &  \\
\midrule
 &  &  &  \\
\midrule
 &  &  &  \\
\midrule
 &  &  &  \\
\midrule
 &  &  &  \\
\midrule
 &  &  &  \\
\midrule
 &  &  &  \\
\midrule
 &  &  &  \\
\midrule
 &  &  &  \\
\midrule
 &  &  &  \\
\midrule
 &  &  &  \\
\midrule
 &  &  &  \\
\midrule
 &  &  &  \\
\midrule
 &  &  &  \\
\midrule
 &  &  &  \\
\midrule
 &  &  &  \\
\midrule
 &  &  &  \\
\midrule
 &  &  &  \\
\midrule
 &  &  &  \\
\midrule
 &  &  &  \\
\midrule
 &  &  &  \\
\midrule
 &  &  &  \\
\midrule
 &  &  &  \\
\midrule
 &  &  &  \\
\midrule
 &  &  &  \\
\midrule
 &  &  &  \\
\midrule
 &  &  &  \\
\midrule
 &  &  &  \\
\midrule
 &  &  &  \\
\midrule
 &  &  &  \\
\midrule
 &  &  &  \\
\midrule
 &  &  &  \\
\midrule
 &  &  &  \\
\midrule
 &  &  &  \\
\midrule
 &  &  &  \\
\midrule
 &  &  &  \\
\midrule
 &  &  &  \\
\midrule
 &  &  &  \\
\midrule
 &  &  &  \\
\midrule
 &  &  &  \\
\midrule
 &  &  &  \\
\midrule
 &  &  &  \\
\midrule
 &  &  &  \\
\midrule
 &  &  &  \\
\midrule
 &  &  &  \\
\midrule
 &  &  &  \\
\midrule
 &  &  &  \\
\midrule
 &  &  &  \\
\midrule
 &  &  &  \\
\midrule
 &  &  &  \\
\midrule
 &  &  &  \\
\midrule
 &  &  &  \\
\midrule
 &  &  &  \\
\midrule
 &  &  &  \\
\midrule
 &  &  &  \\
\midrule
 &  &  &  \\
\midrule
 &  &  &  \\
\midrule
 &  &  &  \\
\midrule
 &  &  &  \\
\midrule
 &  &  &  \\
\midrule
 &  &  &  \\
\midrule
 &  &  &  \\
\midrule
 &  &  &  \\
\midrule
 &  &  &  \\
\midrule
 &  &  &  \\
\midrule
 &  &  &  \\
\midrule
 &  &  &  \\
\midrule
 &  &  &  \\
\midrule
 &  &  &  \\
\midrule
 &  &  &  \\
\midrule
 &  &  &  \\
\midrule
 &  &  &  \\
\midrule
 &  &  &  \\
\midrule
 &  &  &  \\
\midrule
 &  &  &  \\
\midrule
 &  &  &  \\
\midrule
 &  &  &  \\
\midrule
 &  &  &  \\
\midrule
 &  &  &  \\
\midrule
 &  &  &  \\
\midrule
 &  &  &  \\
\midrule
 &  &  &  \\
\midrule
 &  &  &  \\
\midrule
 &  &  &  \\
\midrule
 &  &  &  \\
\midrule
 &  &  &  \\
\midrule
 &  &  &  \\
\midrule
 &  &  &  \\
\midrule
 &  &  &  \\
\midrule
 &  &  &  \\
\midrule
 &  &  &  \\
\midrule
 &  &  &  \\
\midrule
 &  &  &  \\
\midrule
 &  &  &  \\
\midrule
 &  &  &  \\
\midrule
 &  &  &  \\
\midrule
 &  &  &  \\
\midrule
 &  &  &  \\
\midrule
 &  &  &  \\
\midrule
 &  &  &  \\
\midrule
 &  &  &  \\
\midrule
 &  &  &  \\
\midrule
 &  &  &  \\
\midrule
 &  &  &  \\
\midrule
 &  &  &  \\
\midrule
 &  &  &  \\
\midrule
 &  &  &  \\
\midrule
 &  &  &  \\
\midrule
 &  &  &  \\
\midrule
 &  &  &  \\
\midrule
 &  &  &  \\
\midrule
 &  &  &  \\
\midrule
 &  &  &  \\
\midrule
 &  &  &  \\
\midrule
 &  &  &  \\
\midrule
 &  &  &  \\
\midrule
 &  &  &  \\
\midrule
 &  &  &  \\
\midrule
 &  &  &  \\
\midrule
 &  &  &  \\
\midrule
 &  &  &  \\
\midrule
 &  &  &  \\
\midrule
 &  &  &  \\
\midrule
 &  &  &  \\
\midrule
 &  &  &  \\
\midrule
 &  &  &  \\
\midrule
 &  &  &  \\
\midrule
 &  &  &  \\
\midrule
 &  &  &  \\
\midrule
 &  &  &  \\
\midrule
 &  &  &  \\
\midrule
 &  &  &  \\
\midrule
 &  &  &  \\
\midrule
 &  &  &  \\
\midrule
 &  &  &  \\
\midrule
 &  &  &  \\
\midrule
 &  &  &  \\
\midrule
 &  &  &  \\
\midrule
 &  &  &  \\
\midrule
 &  &  &  \\
\midrule
 &  &  &  \\
\midrule
 &  &  &  \\
\midrule
 &  &  &  \\
\midrule
 &  &  &  \\
\midrule
 &  &  &  \\
\midrule
 &  &  &  \\
\midrule
 &  &  &  \\
\midrule
 &  &  &  \\
\midrule
 &  &  &  \\
\midrule
 &  &  &  \\
\midrule
 &  &  &  \\
\midrule
 &  &  &  \\
\midrule
 &  &  &  \\
\midrule
 &  &  &  \\
\midrule
 &  &  &  \\
\midrule
 &  &  &  \\
\midrule
 &  &  &  \\
\midrule
 &  &  &  \\
\midrule
 &  &  &  \\
\midrule
 &  &  &  \\
\midrule
 &  &  &  \\
\midrule
 &  &  &  \\
\midrule
 &  &  &  \\
\midrule
 &  &  &  \\
\midrule
 &  &  &  \\
\midrule
 &  &  &  \\
\midrule
 &  &  &  \\
\midrule
 &  &  &  \\
\midrule
 &  &  &  \\
\midrule
 &  &  &  \\
\midrule
 &  &  &  \\
\midrule
 &  &  &  \\
\midrule
 &  &  &  \\
\midrule
 &  &  &  \\
\midrule
 &  &  &  \\
\midrule
 &  &  &  \\
\midrule
 &  &  &  \\
\midrule
 &  &  &  \\
\midrule
 &  &  &  \\
\midrule
 &  &  &  \\
\midrule
 &  &  &  \\
\midrule
 &  &  &  \\
\midrule
 &  &  &  \\
\midrule
 &  &  &  \\
\midrule
 &  &  &  \\
\midrule
 &  &  &  \\
\midrule
 &  &  &  \\
\midrule
 &  &  &  \\
\midrule
 &  &  &  \\
\midrule
 &  &  &  \\
\midrule
 &  &  &  \\
\midrule
 &  &  &  \\
\midrule
 &  &  &  \\
\midrule
 &  &  &  \\
\midrule
 &  &  &  \\
\midrule
 &  &  &  \\
\midrule
 &  &  &  \\
\midrule
 &  &  &  \\
\midrule
 &  &  &  \\
\midrule
 &  &  &  \\
\midrule
 &  &  &  \\
\midrule
 &  &  &  \\
\midrule
 &  &  &  \\
\midrule
 &  &  &  \\
\midrule
 &  &  &  \\
\midrule
 &  &  &  \\
\midrule
 &  &  &  \\
\midrule
 &  &  &  \\
\midrule
 &  &  &  \\
\midrule
 &  &  &  \\
\midrule
 &  &  &  \\
\midrule
 &  &  &  \\
\midrule
 &  &  &  \\
\midrule
 &  &  &  \\
\midrule
 &  &  &  \\
\midrule
 &  &  &  \\
\midrule
 &  &  &  \\
\midrule
 &  &  &  \\
\midrule
 &  &  &  \\
\midrule
 &  &  &  \\
\midrule
 &  &  &  \\
\midrule
 &  &  &  \\
\midrule
 &  &  &  \\
\midrule
 &  &  &  \\
\midrule
 &  &  &  \\
\midrule
 &  &  &  \\
\midrule
 &  &  &  \\
\midrule
 &  &  &  \\
\midrule
 &  &  &  \\
\midrule
 &  &  &  \\
\midrule
 &  &  &  \\
\midrule
 &  &  &  \\
\midrule
 &  &  &  \\
\midrule
 &  &  &  \\
\midrule
 &  &  &  \\
\midrule
 &  &  &  \\
\midrule
 &  &  &  \\
\midrule
 &  &  &  \\
\midrule
 &  &  &  \\
\midrule
 &  &  &  \\
\midrule
 &  &  &  \\
\midrule
 &  &  &  \\
\midrule
 &  &  &  \\
\midrule
 &  &  &  \\
\midrule
 &  &  &  \\
\midrule
 &  &  &  \\
\midrule
 &  &  &  \\
\midrule
 &  &  &  \\
\midrule
 &  &  &  \\
\midrule
 &  &  &  \\
\midrule
 &  &  &  \\
\midrule
 &  &  &  \\
\midrule
 &  &  &  \\
\midrule
 &  &  &  \\
\midrule
 &  &  &  \\
\midrule
 &  &  &  \\
\midrule
 &  &  &  \\
\midrule
 &  &  &  \\
\midrule
 &  &  &  \\
\midrule
 &  &  &  \\
\midrule
 &  &  &  \\
\midrule
 &  &  &  \\
\midrule
 &  &  &  \\
\midrule
 &  &  &  \\
\midrule
 &  &  &  \\
\midrule
 &  &  &  \\
\midrule
 &  &  &  \\
\midrule
 &  &  &  \\
\midrule
 &  &  &  \\
\midrule
 &  &  &  \\
\midrule
 &  &  &  \\
\midrule
 &  &  &  \\
\midrule
 &  &  &  \\
\midrule
 &  &  &  \\
\midrule
 &  &  &  \\
\midrule
 &  &  &  \\
\midrule
 &  &  &  \\
\midrule
 &  &  &  \\
\midrule
 &  &  &  \\
\midrule
 &  &  &  \\
\midrule
 &  &  &  \\
\midrule
 &  &  &  \\
\midrule
 &  &  &  \\
\midrule
 &  &  &  \\
\midrule
 &  &  &  \\
\midrule
 &  &  &  \\
\midrule
 &  &  &  \\
\midrule
 &  &  &  \\
\midrule
 &  &  &  \\
\midrule
 &  &  &  \\
\midrule
 &  &  &  \\
\midrule
 &  &  &  \\
\midrule
 &  &  &  \\
\midrule
 &  &  &  \\
\midrule
 &  &  &  \\
\midrule
 &  &  &  \\
\midrule
 &  &  &  \\
\midrule
 &  &  &  \\
\midrule
 &  &  &  \\
\midrule
 &  &  &  \\
\midrule
 &  &  &  \\
\midrule
 &  &  &  \\
\midrule
 &  &  &  \\
\midrule
 &  &  &  \\
\midrule
 &  &  &  \\
\midrule
 &  &  &  \\
\midrule
 &  &  &  \\
\midrule
 &  &  &  \\
\midrule
 &  &  &  \\
\midrule
 &  &  &  \\
\midrule
 &  &  &  \\
\midrule
 &  &  &  \\
\midrule
 &  &  &  \\
\midrule
 &  &  &  \\
\midrule
 &  &  &  \\
\midrule
 &  &  &  \\
\midrule
 &  &  &  \\
\midrule
 &  &  &  \\
\midrule
 &  &  &  \\
\midrule
 &  &  &  \\
\midrule
 &  &  &  \\
\midrule
 &  &  &  \\
\midrule
 &  &  &  \\
\midrule
 &  &  &  \\
\midrule
 &  &  &  \\
\midrule
 &  &  &  \\
\midrule
 &  &  &  \\
\midrule
 &  &  &  \\
\midrule
 &  &  &  \\
\midrule
 &  &  &  \\
\midrule
 &  &  &  \\
\midrule
 &  &  &  \\
\midrule
 &  &  &  \\
\midrule
 &  &  &  \\
\midrule
 &  &  &  \\
\midrule
 &  &  &  \\
\midrule
 &  &  &  \\
\midrule
 &  &  &  \\
\midrule
 &  &  &  \\
\midrule
 &  &  &  \\
\midrule
 &  &  &  \\
\midrule
 &  &  &  \\
\midrule
 &  &  &  \\
\midrule
 &  &  &  \\
\midrule
 &  &  &  \\
\midrule
 &  &  &  \\
\midrule
 &  &  &  \\
\midrule
 &  &  &  \\
\midrule
 &  &  &  \\
\midrule
 &  &  &  \\
\midrule
 &  &  &  \\
\midrule
 &  &  &  \\
\midrule
 &  &  &  \\
\midrule
 &  &  &  \\
\midrule
 &  &  &  \\
\midrule
 &  &  &  \\
\midrule
 &  &  &  \\
\midrule
 &  &  &  \\
\midrule
 &  &  &  \\
\midrule
 &  &  &  \\
\midrule
 &  &  &  \\
\midrule
 &  &  &  \\
\midrule
 &  &  &  \\
\midrule
 &  &  &  \\
\midrule
 &  &  &  \\
\midrule
 &  &  &  \\
\midrule
 &  &  &  \\
\midrule
 &  &  &  \\
\midrule
 &  &  &  \\
\midrule
 &  &  &  \\
\midrule
 &  &  &  \\
\midrule
 &  &  &  \\
\midrule
 &  &  &  \\
\midrule
 &  &  &  \\
\midrule
 &  &  &  \\
\midrule
 &  &  &  \\
\midrule
 &  &  &  \\
\midrule
 &  &  &  \\
\midrule
 &  &  &  \\
\midrule
 &  &  &  \\
\midrule
 &  &  &  \\
\midrule
 &  &  &  \\
\midrule
 &  &  &  \\
\midrule
 &  &  &  \\
\midrule
 &  &  &  \\
\midrule
 &  &  &  \\
\midrule
 &  &  &  \\
\midrule
 &  &  &  \\
\midrule
 &  &  &  \\
\midrule
 &  &  &  \\
\midrule
 &  &  &  \\
\midrule
 &  &  &  \\
\midrule
 &  &  &  \\
\midrule
 &  &  &  \\
\midrule
 &  &  &  \\
\midrule
 &  &  &  \\
\midrule
 &  &  &  \\
\midrule
 &  &  &  \\
\midrule
 &  &  &  \\
\midrule
 &  &  &  \\
\midrule
 &  &  &  \\
\midrule
 &  &  &  \\
\midrule
 &  &  &  \\
\midrule
 &  &  &  \\
\midrule
 &  &  &  \\
\midrule
 &  &  &  \\
\midrule
 &  &  &  \\
\midrule
 &  &  &  \\
\midrule
 &  &  &  \\
\midrule
 &  &  &  \\
\midrule
 &  &  &  \\
\midrule
 &  &  &  \\
\midrule
 &  &  &  \\
\midrule
 &  &  &  \\
\midrule
 &  &  &  \\
\midrule
 &  &  &  \\
\midrule
 &  &  &  \\
\midrule
 &  &  &  \\
\midrule
 &  &  &  \\
\midrule
 &  &  &  \\
\midrule
 &  &  &  \\
\midrule
 &  &  &  \\
\midrule
 &  &  &  \\
\midrule
 &  &  &  \\
\midrule
 &  &  &  \\
\midrule
 &  &  &  \\
\midrule
 &  &  &  \\
\midrule
 &  &  &  \\
\midrule
 &  &  &  \\
\midrule
 &  &  &  \\
\midrule
 &  &  &  \\
\midrule
 &  &  &  \\
\midrule
 &  &  &  \\
\midrule
 &  &  &  \\
\midrule
 &  &  &  \\
\midrule
 &  &  &  \\
\midrule
 &  &  &  \\
\midrule
 &  &  &  \\
\midrule
 &  &  &  \\
\midrule
 &  &  &  \\
\midrule
 &  &  &  \\
\midrule
 &  &  &  \\
\midrule
 &  &  &  \\
\midrule
 &  &  &  \\
\midrule
 &  &  &  \\
\midrule
 &  &  &  \\
\midrule
 &  &  &  \\
\midrule
 &  &  &  \\
\midrule
 &  &  &  \\
\midrule
 &  &  &  \\
\midrule
 &  &  &  \\
\midrule
 &  &  &  \\
\midrule
 &  &  &  \\
\midrule
 &  &  &  \\
\midrule
 &  &  &  \\
\midrule
 &  &  &  \\
\midrule
 &  &  &  \\
\midrule
 &  &  &  \\
\midrule
 &  &  &  \\
\midrule
 &  &  &  \\
\midrule
 &  &  &  \\
\midrule
 &  &  &  \\
\midrule
 &  &  &  \\
\midrule
 &  &  &  \\
\midrule
 &  &  &  \\
\midrule
 &  &  &  \\
\midrule
 &  &  &  \\
\midrule
 &  &  &  \\
\midrule
 &  &  &  \\
\midrule
 &  &  &  \\
\midrule
 &  &  &  \\
\midrule
 &  &  &  \\
\midrule
 &  &  &  \\
\midrule
 &  &  &  \\
\midrule
 &  &  &  \\
\midrule
 &  &  &  \\
\midrule
 &  &  &  \\
\midrule
 &  &  &  \\
\midrule
 &  &  &  \\
\midrule
 &  &  &  \\
\midrule
 &  &  &  \\
\midrule
 &  &  &  \\
\midrule
 &  &  &  \\
\midrule
 &  &  &  \\
\midrule
 &  &  &  \\
\midrule
 &  &  &  \\
\midrule
 &  &  &  \\
\midrule
 &  &  &  \\
\midrule
 &  &  &  \\
\midrule
 &  &  &  \\
\midrule
 &  &  &  \\
\midrule
 &  &  &  \\
\midrule
 &  &  &  \\
\midrule
 &  &  &  \\
\midrule
 &  &  &  \\
\midrule
 &  &  &  \\
\midrule
 &  &  &  \\
\midrule
 &  &  &  \\
\midrule
 &  &  &  \\
\midrule
 &  &  &  \\
\midrule
 &  &  &  \\
\midrule
 &  &  &  \\
\midrule
 &  &  &  \\
\midrule
 &  &  &  \\
\midrule
 &  &  &  \\
\midrule
 &  &  &  \\
\midrule
 &  &  &  \\
\midrule
 &  &  &  \\
\midrule
 &  &  &  \\
\midrule
 &  &  &  \\
\midrule
 &  &  &  \\
\midrule
 &  &  &  \\
\midrule
 &  &  &  \\
\midrule
 &  &  &  \\
\midrule
 &  &  &  \\
\midrule
 &  &  &  \\
\midrule
 &  &  &  \\
\midrule
 &  &  &  \\
\midrule
 &  &  &  \\
\midrule
 &  &  &  \\
\midrule
 &  &  &  \\
\midrule
 &  &  &  \\
\midrule
 &  &  &  \\
\midrule
 &  &  &  \\
\midrule
 &  &  &  \\
\midrule
 &  &  &  \\
\midrule
 &  &  &  \\
\midrule
 &  &  &  \\
\midrule
 &  &  &  \\
\midrule
 &  &  &  \\
\midrule
 &  &  &  \\
\midrule
 &  &  &  \\
\midrule
 &  &  &  \\
\midrule
 &  &  &  \\
\midrule
 &  &  &  \\
\midrule
 &  &  &  \\
\midrule
 &  &  &  \\
\midrule
 &  &  &  \\
\midrule
 &  &  &  \\
\midrule
 &  &  &  \\
\midrule
 &  &  &  \\
\midrule
 &  &  &  \\
\midrule
 &  &  &  \\
\midrule
 &  &  &  \\
\midrule
 &  &  &  \\
\midrule
 &  &  &  \\
\midrule
 &  &  &  \\
\midrule
 &  &  &  \\
\midrule
 &  &  &  \\
\midrule
 &  &  &  \\
\midrule
 &  &  &  \\
\midrule
 &  &  &  \\
\midrule
 &  &  &  \\
\midrule
 &  &  &  \\
\midrule
 &  &  &  \\
\midrule
 &  &  &  \\
\midrule
 &  &  &  \\
\midrule
 &  &  &  \\
\midrule
 &  &  &  \\
\midrule
 &  &  &  \\
\midrule
 &  &  &  \\
\midrule
 &  &  &  \\
\midrule
 &  &  &  \\
\midrule
 &  &  &  \\
\midrule
 &  &  &  \\
\midrule
 &  &  &  \\
\midrule
 &  &  &  \\
\midrule
 &  &  &  \\
\midrule
 &  &  &  \\
\midrule
 &  &  &  \\
\midrule
 &  &  &  \\
\midrule
 &  &  &  \\
\midrule
 &  &  &  \\
\midrule
 &  &  &  \\
\midrule
 &  &  &  \\
\midrule
 &  &  &  \\
\midrule
 &  &  &  \\
\midrule
 &  &  &  \\
\midrule
 &  &  &  \\
\midrule
 &  &  &  \\
\midrule
 &  &  &  \\
\midrule
 &  &  &  \\
\midrule
 &  &  &  \\
\midrule
 &  &  &  \\
\midrule
 &  &  &  \\
\midrule
 &  &  &  \\
\midrule
 &  &  &  \\
\midrule
 &  &  &  \\
\midrule
 &  &  &  \\
\midrule
 &  &  &  \\
\midrule
 &  &  &  \\
\midrule
 &  &  &  \\
\midrule
 &  &  &  \\
\midrule
 &  &  &  \\
\midrule
 &  &  &  \\
\midrule
 &  &  &  \\
\midrule
 &  &  &  \\
\midrule
 &  &  &  \\
\midrule
 &  &  &  \\
\midrule
 &  &  &  \\
\midrule
 &  &  &  \\
\midrule
 &  &  &  \\
\midrule
 &  &  &  \\
\midrule
 &  &  &  \\
\midrule
 &  &  &  \\
\midrule
 &  &  &  \\
\midrule
 &  &  &  \\
\midrule
 &  &  &  \\
\midrule
 &  &  &  \\
\midrule
 &  &  &  \\
\midrule
 &  &  &  \\
\midrule
 &  &  &  \\
\midrule
 &  &  &  \\
\midrule
 &  &  &  \\
\midrule
 &  &  &  \\
\midrule
 &  &  &  \\
\midrule
 &  &  &  \\
\midrule
 &  &  &  \\
\midrule
 &  &  &  \\
\midrule
 &  &  &  \\
\midrule
 &  &  &  \\
\midrule
 &  &  &  \\
\midrule
 &  &  &  \\
\midrule
 &  &  &  \\
\midrule
 &  &  &  \\
\midrule
 &  &  &  \\
\midrule
 &  &  &  \\
\midrule
 &  &  &  \\
\midrule
 &  &  &  \\
\midrule
 &  &  &  \\
\midrule
 &  &  &  \\
\midrule
 &  &  &  \\
\midrule
 &  &  &  \\
\midrule
 &  &  &  \\
\midrule
 &  &  &  \\
\midrule
 &  &  &  \\
\midrule
 &  &  &  \\
\midrule
 &  &  &  \\
\midrule
 &  &  &  \\
\midrule
 &  &  &  \\
\midrule
 &  &  &  \\
\midrule
 &  &  &  \\
\midrule
 &  &  &  \\
\midrule
 &  &  &  \\
\midrule
 &  &  &  \\
\midrule
 &  &  &  \\
\midrule
 &  &  &  \\
\midrule
 &  &  &  \\
\midrule
 &  &  &  \\
\midrule
 &  &  &  \\
\midrule
 &  &  &  \\
\midrule
 &  &  &  \\
\midrule
 &  &  &  \\
\midrule
 &  &  &  \\
\midrule
 &  &  &  \\
\midrule
 &  &  &  \\
\midrule
 &  &  &  \\
\midrule
 &  &  &  \\
\midrule
 &  &  &  \\
\midrule
 &  &  &  \\
\midrule
 &  &  &  \\
\midrule
 &  &  &  \\
\midrule
 &  &  &  \\
\midrule
 &  &  &  \\
\midrule
 &  &  &  \\
\midrule
 &  &  &  \\
\midrule
 &  &  &  \\
\midrule
 &  &  &  \\
\midrule
 &  &  &  \\
\midrule
 &  &  &  \\
\midrule
 &  &  &  \\
\midrule
 &  &  &  \\
\midrule
 &  &  &  \\
\midrule
 &  &  &  \\
\midrule
 &  &  &  \\
\midrule
 &  &  &  \\
\midrule
 &  &  &  \\
\midrule
 &  &  &  \\
\midrule
 &  &  &  \\
\midrule
 &  &  &  \\
\midrule
 &  &  &  \\
\midrule
 &  &  &  \\
\midrule
 &  &  &  \\
\midrule
 &  &  &  \\
\midrule
 &  &  &  \\
\midrule
 &  &  &  \\
\midrule
 &  &  &  \\
\midrule
 &  &  &  \\
\midrule
 &  &  &  \\
\midrule
 &  &  &  \\
\midrule
 &  &  &  \\
\midrule
 &  &  &  \\
\midrule
 &  &  &  \\
\midrule
 &  &  &  \\
\midrule
 &  &  &  \\
\midrule
 &  &  &  \\
\midrule
 &  &  &  \\
\midrule
 &  &  &  \\
\midrule
 &  &  &  \\
\midrule
 &  &  &  \\
\midrule
 &  &  &  \\
\midrule
 &  &  &  \\
\midrule
 &  &  &  \\
\midrule
 &  &  &  \\
\midrule
 &  &  &  \\
\midrule
 &  &  &  \\
\midrule
 &  &  &  \\
\midrule
 &  &  &  \\
\midrule
 &  &  &  \\
\midrule
 &  &  &  \\
\midrule
 &  &  &  \\
\midrule
 &  &  &  \\
\midrule
 &  &  &  \\
\midrule
 &  &  &  \\
\midrule
 &  &  &  \\
\midrule
 &  &  &  \\
\midrule
 &  &  &  \\
\midrule
 &  &  &  \\
\midrule
 &  &  &  \\
\midrule
 &  &  &  \\
\midrule
 &  &  &  \\
\midrule
 &  &  &  \\
\midrule
 &  &  &  \\
\midrule
 &  &  &  \\
\midrule
 &  &  &  \\
\midrule
 &  &  &  \\
\midrule
 &  &  &  \\
\midrule
 &  &  &  \\
\midrule
 &  &  &  \\
\midrule
 &  &  &  \\
\midrule
 &  &  &  \\
\midrule
 &  &  &  \\
\midrule
 &  &  &  \\
\midrule
 &  &  &  \\
\midrule
 &  &  &  \\
\midrule
 &  &  &  \\
\midrule
 &  &  &  \\
\midrule
 &  &  &  \\
\midrule
 &  &  &  \\
\midrule
 &  &  &  \\
\midrule
 &  &  &  \\
\midrule
 &  &  &  \\
\midrule
 &  &  &  \\
\midrule
 &  &  &  \\
\midrule
 &  &  &  \\
\midrule
 &  &  &  \\
\midrule
 &  &  &  \\
\midrule
 &  &  &  \\
\midrule
 &  &  &  \\
\midrule
 &  &  &  \\
\midrule
 &  &  &  \\
\midrule
 &  &  &  \\
\midrule
 &  &  &  \\
\midrule
 &  &  &  \\
\midrule
 &  &  &  \\
\midrule
 &  &  &  \\
\midrule
 &  &  &  \\
\midrule
 &  &  &  \\
\midrule
 &  &  &  \\
\midrule
 &  &  &  \\
\midrule
 &  &  &  \\
\midrule
 &  &  &  \\
\midrule
 &  &  &  \\
\midrule
 &  &  &  \\
\midrule
 &  &  &  \\
\midrule
 &  &  &  \\
\midrule
 &  &  &  \\
\midrule
 &  &  &  \\
\midrule
 &  &  &  \\
\midrule
 &  &  &  \\
\midrule
 &  &  &  \\
\midrule
 &  &  &  \\
\midrule
 &  &  &  \\
\midrule
 &  &  &  \\
\midrule
 &  &  &  \\
\midrule
 &  &  &  \\
\midrule
 &  &  &  \\
\midrule
 &  &  &  \\
\midrule
 &  &  &  \\
\midrule
 &  &  &  \\
\midrule
 &  &  &  \\
\midrule
 &  &  &  \\
\midrule
 &  &  &  \\
\midrule
 &  &  &  \\
\midrule
 &  &  &  \\
\midrule
 &  &  &  \\
\midrule
 &  &  &  \\
\midrule
 &  &  &  \\
\midrule
 &  &  &  \\
\midrule
 &  &  &  \\
\midrule
 &  &  &  \\
\midrule
 &  &  &  \\
\midrule
 &  &  &  \\
\midrule
 &  &  &  \\
\midrule
 &  &  &  \\
\midrule
 &  &  &  \\
\midrule
 &  &  &  \\
\midrule
 &  &  &  \\
\midrule
 &  &  &  \\
\midrule
 &  &  &  \\
\midrule
 &  &  &  \\
\midrule
 &  &  &  \\
\midrule
 &  &  &  \\
\midrule
 &  &  &  \\
\midrule
 &  &  &  \\
\midrule
 &  &  &  \\
\midrule
 &  &  &  \\
\midrule
 &  &  &  \\
\midrule
 &  &  &  \\
\midrule
 &  &  &  \\
\midrule
 &  &  &  \\
\midrule
 &  &  &  \\
\midrule
 &  &  &  \\
\midrule
 &  &  &  \\
\midrule
 &  &  &  \\
\midrule
 &  &  &  \\
\midrule
 &  &  &  \\
\midrule
 &  &  &  \\
\midrule
 &  &  &  \\
\midrule
 &  &  &  \\
\midrule
 &  &  &  \\
\midrule
 &  &  &  \\
\midrule
 &  &  &  \\
\midrule
 &  &  &  \\
\midrule
 &  &  &  \\
\midrule
 &  &  &  \\
\midrule
 &  &  &  \\
\midrule
 &  &  &  \\
\midrule
 &  &  &  \\
\midrule
 &  &  &  \\
\midrule
\bottomrule
\end{tabularx}


\section{Cómo se calcula una máscara}\label{sec:calculo-mascara}
Una máscara es un número de 32 bits. Su bit número $n$ vale 1 si el byte $n$ del bloque participa. Como los bloques tienen 28 bytes, solo se usan los 28 bits bajos (por eso el valor «todos» es \cod{16\#0FFFFFFF}, y no \cod{16\#FFFFFFFF}).

\subsection*{Método en cuatro pasos}
\begin{enumerate}
\item \textbf{Lista} los bytes que tu estación \emph{escribe}.
\item \textbf{Lista} los bytes que tu estación \emph{lee}, y añade los que escribes (regla de oro 1, más abajo). Esa es la lista de lectura.
\item Para cada lista, \textbf{busca el valor de cada byte} en la tabla~\ref{tab:valores-byte} y \textbf{súmalos por grupos de cuatro bytes} (cada grupo es un dígito hexadecimal).
\item \textbf{Escribe el resultado} con ocho dígitos: \cod{16\#} seguido de los siete grupos de la derecha a la izquierda.
\end{enumerate}

\begin{table}[ht]
\centering
\footnotesize\setlength{\tabcolsep}{2.0pt}
\begin{tabular}{l|cccc|cccc|cccc|cccc|cccc|cccc|cccc}
\toprule
\textbf{Dígito} (de dcha. a izq.) & \multicolumn{4}{c|}{1.º} & \multicolumn{4}{c|}{2.º} & \multicolumn{4}{c|}{3.º} & \multicolumn{4}{c|}{4.º} & \multicolumn{4}{c|}{5.º} & \multicolumn{4}{c|}{6.º} & \multicolumn{4}{c}{7.º} \\
\textbf{Bytes} & 0&1&2&3 & 4&5&6&7 & 8&9&10&11 & 12&13&14&15 & 16&17&18&19 & 20&21&22&23 & 24&25&26&27 \\
\textbf{Valor} & 1&2&4&8 & 1&2&4&8 & 1&2&4&8 & 1&2&4&8 & 1&2&4&8 & 1&2&4&8 & 1&2&4&8 \\
\bottomrule
\end{tabular}
\caption{Valor de cada byte dentro de su dígito hexadecimal. Los bytes de cada grupo de cuatro suman en su dígito: 1, 2, 4 y 8.}
\label{tab:valores-byte}
\end{table}

\begin{ejemplo}[title={Ejemplo: la máscara de lectura de Press}]
Press \emph{lee} los bytes 3, 12, 13, 14 y 15, y \emph{escribe} los bytes 4 y 16. La lista de lectura es la unión: 3, 4, 12, 13, 14, 15, 16. Ahora, dígito a dígito (de la derecha):
\begin{center}
\small
\begin{tabular}{llcl}
\toprule
Dígito & Bytes de este grupo en la lista & Suma & Resultado \\
\midrule
1.º (bytes 0--3) & 3 & 8 & \cod{8} \\
2.º (bytes 4--7) & 4 & 1 & \cod{1} \\
3.º (bytes 8--11) & ninguno & 0 & \cod{0} \\
4.º (bytes 12--15) & 12, 13, 14, 15 & 1+2+4+8 = 15 & \cod{F} \\
5.º (bytes 16--19) & 16 & 1 & \cod{1} \\
6.º y 7.º & ninguno & 0 & \cod{0}\ \cod{0} \\
\bottomrule
\end{tabular}
\end{center}
Leyendo de izquierda a derecha, quedan los dígitos 0, 0, 1, F, 0, 1, 8: la máscara es \cod{16\#0001F018}. La de escritura (bytes 4 y 16) sale igual: \cod{16\#00010010}.
\end{ejemplo}

\begin{consejo}
Para comprobar una máscara ya hecha, haz el camino contrario: pásala a binario dígito a dígito y mira qué bits valen 1. Si salen los bytes que esperabas, está bien. El simulador del grupo hace exactamente esta comprobación con las máscaras de las tablas siguientes.
\end{consejo}

\section{Las máscaras de las cuatro estaciones}
Estas son las máscaras que salen de aplicar el método a la lógica del Excel. Son los valores con los que se han hecho las pruebas del capítulo~\ref{cap:pruebas}.

\begin{tabularx}{\textwidth}{>{\bfseries}l p{3.1cm} c c X}
\toprule
Estación & Bytes que usa & \texttt{Mascara\_lectura} & \texttt{Mascara\_escritura} & Qué escribe \\
\midrule
Cámara & 0, 1, 2, 3, 4, 5, 6, 8, 9, 10, 11, 20 & \texttt{16\#00100F7F} & \texttt{16\#00000E7F} & bytes 0 a 6 (los 0 a 3 solo si el tag está vacío), 9, 10 y 11 \\
BackCover & 3, 4, 12, 23 & \texttt{16\#00801018} & \texttt{16\#00001010} & byte 4 (bit 1) y byte 12 \\
Press & 3, 4, 12, 13, 14, 15, 16 & \texttt{16\#0001F018} & \texttt{16\#00010010} & byte 4 (bit 2) y byte 16 \\
Out & 5, 9, 17, 21, 22 & \texttt{16\#00620220} & \texttt{16\#0FFFFFFF} & los 28 bytes: deja el bloque a cero \\
\bottomrule
\end{tabularx}


\begin{ojo}
La estación \textbf{Out} es distinta a propósito. Lee solo los bytes que necesita para decidir, y al final \textbf{escribe los 28 bytes a cero}: según el Excel, «Out pone el bloque a cero al depositar la pieza, para que el palet empiece limpio la siguiente vuelta». Al escribir ceros, la imagen que escribe tiene que estar entera a cero (su aplicación la pone a cero en la etapa de fin) y por eso su máscara de escritura es \cod{16\#0FFFFFFF}. Esto deja también a cero los bytes 20 a 27, que son de otros grupos: cuando el palet sale de Out ya no lleva ningún dato de nadie.
\end{ojo}

\section{Las reglas de oro}\label{sec:reglas-oro}
\begin{enumerate}[leftmargin=1.2cm,label=\textbf{O\arabic*.}]
\item \textbf{Todo byte que escribes, también lo lees.} Si un byte está en \cod{Mascara\_escritura} y no en \cod{Mascara\_lectura}, la imagen lo tiene a cero cuando tu aplicación lo recibe. Si no lo rellenas entero tú, escribirás ceros encima del valor del tag. Excepción: bytes que tu aplicación \emph{siempre} rellena entera, como hace Out al poner a cero el bloque.
\item \textbf{Escribe solo los bytes que son tuyos.} Mira la columna «Lo escribe» del mapa. Si tu estación tiene que escribir en un byte que es de otra, habla antes con su dueño.
\item \textbf{Las banderas se completan con OR, no con una asignación.} En el byte 4 (estaciones completadas) cada estación pone \emph{su} bit sin tocar los demás: \texttt{byte4 := byte4 OR mi\_bit}, nunca \texttt{byte4 := valor}. Por eso el byte 4 tiene que estar en tu lectura. Lo mismo para el byte 12 (estado de BackCover) y el 9 (resultado de la cámara).
\item \textbf{Valores 0 = «sin dato».} El mapa define para cada byte qué significa que valga 0 (columna «Si vale 0»). Una estación que se encuentre un 0 debe actuar con ese significado (por ejemplo, 0 golpes en la prensa significa «1 golpe»).
\item \textbf{La misma dirección para leer y escribir.} \cod{Direccion\_lectura} y \cod{Direccion\_escritura} deben ser iguales. Si fueran distintas, los bytes que no son tuyos se escribirían en otra zona del tag con datos que no corresponden.
\item \textbf{Los bytes 20 a 27 son una propuesta.} Hasta que los otros grupos acepten el reparto, tu estación debe funcionar sin ellos usando el valor «si vale 0».
\item \textbf{Los números de varios bytes se leen y escriben juntos} (bytes 14--15).
\end{enumerate}

\begin{center}
\small
\begin{tabularx}{\textwidth}{>{\bfseries}p{1.0cm}X}
\toprule
Regla & Qué se nota si la incumples \\
\midrule
O1 & Un byte tuyo aparece a cero en el tag después de pasar por tu estación. \\
O2 & Otra estación encuentra sus datos cambiados sin que ella los haya tocado. \\
O3 & Los bits de otras estaciones desaparecen de «completadas»; la estación siguiente cree que no se ha hecho su paso previo. \\
O4 & Una estación se comporta mal con un palet nuevo o con el tag a cero (por ejemplo, 0 golpes en la prensa). \\
O5 & Los bytes ajenos del tag se estropean sin motivo aparente. \\
\bottomrule
\end{tabularx}
\end{center}

\section{La dirección del bloque en el tag}
El tag tiene 108 bytes y la librería usa 28. La \textbf{dirección de inicio} (\cod{Direccion\_lectura} y \cod{Direccion\_escritura}) dice dónde empieza el bloque. Por defecto es 0, es decir, el bloque ocupa los bytes 0 a 27 del tag. Si en algún momento el profesor o el grupo asignan otro punto de inicio (la hoja de cálculo lo anota en la celda «Dirección de inicio del bloque en el tag»), basta con poner el mismo valor en los dos pines en \cod{Main}, en todas las estaciones.

\begin{ojo}
El manual de la CP-Lab dice en un apartado que los accesos a un tag se hacen «en bloques de 24 bytes», y el resto del manual y las señales de E/S del lector (\cod{\%E14.0}, matriz de 28 bytes) hablan de 28. La librería usa 28, como el UDT del lector y el Excel. Si en la maqueta real se observara que el lector solo transfiere 24, habría que revisarlo (apéndice~\ref{ap:abiertos}).
\end{ojo}

\section{Un palet recorriendo la línea}
Para ver todo junto, la tabla~\ref{tab:recorrido} muestra el bloque de un palet al pasar por las cuatro estaciones, en este orden: Cámara, BackCover, Press y Out. Se ha obtenido ejecutando el código real de la librería y de los ejemplos de aplicación en el simulador, con el palet saliendo de la HMI con ID 7, receta 1, todas las estaciones por hacer, dos golpes de prensa de 2000\,ms. Las celdas con fondo naranja son los bytes que cambian en cada paso.

\begin{table}[ht]
\centering
\footnotesize
\begin{tabular}{r l r r r r r}
\toprule
Byte & Campo & \textbf{Inicio} & \textbf{Cámara} & \textbf{BackCover} & \textbf{Press} & \textbf{Out} \\
\midrule
0 & Marca de bloque válido & \texttt{A5} & \texttt{A5} & \texttt{A5} & \texttt{A5} & \cellcolor{naranja!25}\texttt{00} \\
1 & ID de palet & \texttt{07} & \texttt{07} & \texttt{07} & \texttt{07} & \cellcolor{naranja!25}\texttt{00} \\
2 & Receta & \texttt{01} & \texttt{01} & \texttt{01} & \texttt{01} & \cellcolor{naranja!25}\texttt{00} \\
3 & Estaciones a realizar & \texttt{0F} & \texttt{0F} & \texttt{0F} & \texttt{0F} & \cellcolor{naranja!25}\texttt{00} \\
4 & Estaciones completadas & \texttt{00} & \cellcolor{naranja!25}\texttt{01} & \cellcolor{naranja!25}\texttt{03} & \cellcolor{naranja!25}\texttt{07} & \cellcolor{naranja!25}\texttt{00} \\
5 & Calidad global & \texttt{00} & \cellcolor{naranja!25}\texttt{01} & \texttt{01} & \texttt{01} & \cellcolor{naranja!25}\texttt{00} \\
9 & Resultado de la cámara & \texttt{00} & \cellcolor{naranja!25}\texttt{07} & \texttt{07} & \texttt{07} & \cellcolor{naranja!25}\texttt{00} \\
10 & ImageID & \texttt{00} & \cellcolor{naranja!25}\texttt{2A} & \texttt{2A} & \texttt{2A} & \cellcolor{naranja!25}\texttt{00} \\
11 & Intentos de inspección & \texttt{00} & \cellcolor{naranja!25}\texttt{01} & \texttt{01} & \texttt{01} & \cellcolor{naranja!25}\texttt{00} \\
12 & Estado BackCover & \texttt{00} & \texttt{00} & \cellcolor{naranja!25}\texttt{03} & \texttt{03} & \cellcolor{naranja!25}\texttt{00} \\
13 & Nº de golpes & \texttt{02} & \texttt{02} & \texttt{02} & \texttt{02} & \cellcolor{naranja!25}\texttt{00} \\
14 & Tiempo de prensado (alto) & \texttt{07} & \texttt{07} & \texttt{07} & \texttt{07} & \cellcolor{naranja!25}\texttt{00} \\
15 & Tiempo de prensado (bajo) & \texttt{D0} & \texttt{D0} & \texttt{D0} & \texttt{D0} & \cellcolor{naranja!25}\texttt{00} \\
16 & Golpes realizados & \texttt{00} & \texttt{00} & \texttt{00} & \cellcolor{naranja!25}\texttt{02} & \cellcolor{naranja!25}\texttt{00} \\
\bottomrule
\end{tabular}

\caption{El bloque del tag (en hexadecimal) antes y después de cada estación.}
\label{tab:recorrido}
\end{table}

Se ve exactamente lo que se espera:
\begin{itemize}
\item \textbf{Cámara} marca su bit en el byte 4 (\cod{01}), pone la calidad (\cod{01} = OK), el resultado (\cod{07} = inspeccionado, OK y bien orientado), el ImageID y un intento.
\item \textbf{BackCover} añade su bit (el byte 4 pasa a \cod{03}) y marca la tapa colocada en el byte 12 (\cod{03} = requerida y colocada).
\item \textbf{Press} añade su bit (\cod{07}) y anota los dos golpes realizados en el byte 16.
\item \textbf{Out} deja el bloque entero a cero, listo para el siguiente palet.
\end{itemize}
Nótese que ninguna estación cambia bytes que no son suyos: el resto de la tabla no se mueve.

\section{Cómo cambiar el mapa}
Si el grupo necesita un campo nuevo o cambiar el significado de uno:
\begin{enumerate}
\item Se acuerda entre todos y se actualiza el Excel (lo guarda siempre quien mantiene el mapa; conviene que lleve la fecha en el nombre).
\item A cada estación a la que afecta se le cambia sus constantes \cod{IDX\_*} y, si procede, sus máscaras. \textbf{La librería no cambia}: no sabe qué significa cada byte.
\item Se vuelven a comprobar los pasos de aceptación del capítulo~\ref{cap:pruebas} en las estaciones afectadas.
\end{enumerate}
La librería solo tendría que cambiar si el número de bytes del bloque dejara de ser 28.
